\onecolumn
\section{Evaluation Data Details}
\label{app:data}

Table~\ref{tab:evaluation-data} summarizes the evaluation groups. The
target and neighboring sets use EMBER's original validation/test division. We
rewrote the target and neighboring contexts as declarative sentence stems while
retaining the original correct answer and three distractors; the SciQ items
were scored in their original interrogative form, which is how the instrument
was validated (Section~\ref{sec:evaluation}). We drew the shared SciQ set from
its validation split with seed \texttt{20260923}; the first 50 retained items
form the selection set and the next 50 form the test set.

\begin{center}
  \centering
  \scriptsize
  \begin{tabular}{llrr}
    \toprule
    Group & Source and format & Selection & Test \\
    \midrule
    Target & EMBER, declarative stem & 50 & 50 \\
    Neighboring & EMBER, declarative stem & 50 & 50 \\
    SciQ & Validation, original question & 50 & 50 \\
    \bottomrule
  \end{tabular}
  \captionof{table}{Evaluation-set construction for each concept. The same SciQ items
  are shared across concepts; target and neighboring items are concept
  specific.}
  \label{tab:evaluation-data}
\end{center}

We verified that selection and test sets share neither question IDs nor
identical question--answer pairs. Two of the 50 SciQ selection items are
natively declarative fill-in-the-blank prompts rather than questions. Each
model for a concept was scored with the same frozen strings and option
order. For example,
``Which river is associated with the founding of Rome?'' was presented as
``The river associated with the founding of Rome is'', and the four candidate
answers included the correct continuation ``The Tiber River''.

For likelihood-ranked accuracy, each option was appended to its context and
scored over option tokens only. For correct-answer NLL and KL, the completion
was stored with one leading space, tokenized with the model tokenizer, and
only answer positions were included. Per-question NLL is the mean over answer
tokens. KL is calculated over all 100,352 vocabulary entries at every answer
position and then averaged across positions. Reported group means and sample
standard deviations are computed across the 50 questions. These details
implement the equations in Section~\ref{sec:evaluation}.
